\begin{codeblock}
template<class T>
class @\libglobal{mixer}@ {
public:
  // \ref{paper.mix}, mixing
  int mix(tag_t, const T& part) const
    requires requires { tag; } && (!@\exposidnc{is-void-v}@<T>) && (alpha_policy<T> == 0) &&
             requires(const blend_policy& policy) { policy.strength(); };
};
\end{codeblock}

\rSec3[paper.mix]{Mixing}

\indexlibrarymember{mix}{mixer}%
\begin{itemdecl}
int mix(tag_t, const T& part) const;
\end{itemdecl}

\begin{itemdescr}
\pnum
\constraints
\begin{itemize}
\item \tcode{requires { tag; }} is \tcode{true},
\item \tcode{\exposid{is-void-v}<T>} is \tcode{false},
\item \tcode{alpha_policy<T> == 0} is \tcode{true},
\item \tcode{requires(const blend_policy& policy) { policy.strength(); }} is \tcode{true}.
\end{itemize}

\pnum
\returns
The blended value.
\end{itemdescr}

\indexlibraryglobal{alpha_policy}%
\begin{itemdecl}
template<class T> inline constexpr int alpha_policy = 0;
\end{itemdecl}

\begin{itemdescr}
\pnum
\remarks
This variable template is the lookup point for the alpha policy of a type. A program may specialize it.
\end{itemdescr}

\begin{codeblock}
class @\libglobal{blend_policy}@ {
public:
  // \ref{paper.blend}, observers
  int strength() const;
};
\end{codeblock}

\rSec3[paper.blend]{Blend policy}

\indexlibrarymember{strength}{blend_policy}%
\begin{itemdecl}
int strength() const;
\end{itemdecl}

\begin{itemdescr}
\pnum
\returns
The configured strength.
\end{itemdescr}
